\section*{Ethics Statement}
\label{sec:ethics}

All texts analyzed in this work are LLM-generated from publicly available Wikipedia seeds; no personal or sensitive data is involved.
Our method estimates the generation depth of a text passage (how many recursive generation steps it has undergone), not the identity of its author.
Individual authorship attribution is outside both the scope and the capability of this approach; privacy implications are minimal.
We acknowledge that depth estimation tools could in principle be used to selectively suppress machine-generated text in contexts where such content is legitimate (e.g., accessibility tools, assistive writing).
We encourage deployment in data governance pipelines where provenance transparency, rather than content removal, is the objective.
All code, data, and trained models will be released upon publication.

\paragraph{Licenses.}
All generator models are publicly available under permissive licenses (Apache 2.0 for Pythia, OLMo, and Mistral; Qwen License for Qwen3; Llama Community License for Llama 3.2).
Wikipedia seed passages are distributed under CC BY-SA 3.0.

\paragraph{Compute budget.}
All experiments were conducted on 4$\times$ NVIDIA RTX 5090 GPUs (32\,GB each), totaling approximately 62 GPU-hours.

\paragraph{Software.}
Experiments used Python 3.12.3, PyTorch 2.10.0 (CUDA 12.8), Transformers 4.57.6, XGBoost 3.2.0, scikit-learn 1.8.0, NumPy 2.2.6, pandas 3.0.2, and SciPy 1.17.1.
